\documentclass[UTF8]{ctexart}
\usepackage[a4paper,margin=2.6cm]{geometry}
\begin{document}

\section*{第八部分：带前瞻的 Dijkstra 算法}

第七节已经证明，使用满足 working-set bound 的堆时，普通 Dijkstra 在时间上可以达到 \(O(m+\log D)\)，也就是匹配第六节的时间下界。但是在比较次数上，普通 Dijkstra 的上界是 \(O(F+\log D)\)，而第六节给出的比较下界是 \(\Omega(F-n+1+\log D)\)。两者之间差了一个 \(\Theta(n)\) 的项。当 \(\log D\) 比 \(n\) 小很多时，这个差距就不能忽略。

第八节的目标是修改 Dijkstra，让比较次数也降到

\[
O(F-n+1+\log D),
\]

从而匹配比较下界。这个修改版本叫做 Dijkstra with lookahead。

\subsection*{8.1 瓶颈点}

这一节重新使用第三节中的瓶颈点。顶点 \(v\) 的层数是从 \(s\) 到 \(v\) 的路径上最少经过的顶点数；如果某一层只有一个顶点，那么这个顶点就是瓶颈点。

Lemma 8.1 说明：如果图中有 \(b\) 个瓶颈点，那么

\[
\log D \ge (n-b)/2.
\]

这个结论的意义是：如果非瓶颈点很多，那么 \(D\) 一定不小，普通 Dijkstra 的比较次数已经足够接近最优。因此，真正麻烦的是另一种情况：图里大部分顶点都是瓶颈点。

接下来作者利用瓶颈点的两个性质。

Lemma 8.2 说明：如果 \(v\) 是瓶颈点，那么它支配所有更高层的顶点。也就是说，从 \(s\) 到更高层顶点的每条路径都必须经过 \(v\)。如果有一条弧 \(wv\) 从更高层指回瓶颈点 \(v\)，那么这条弧是无用弧。

Lemma 8.3 说明：如果一个瓶颈点 \(v\) 是 unmarked，而且它不是最高层顶点，那么存在一条弧 \(vw\)，使得

\[
d^*(w)=d^*(v)+c(vw).
\]

这里 marked / unmarked 的定义是：如果瓶颈点的下一层至少有两个顶点，就把它标记为 marked；否则它是 unmarked。直观上，连续的 unmarked 瓶颈点形成一条很窄的链。一旦知道链上第一个瓶颈点的真实距离，就可以沿着层数顺序继续算出后面瓶颈点的真实距离。

这就是 lookahead 的核心想法：瓶颈点不必像普通顶点一样都放进堆里等 \texttt{delete-min}，而是可以单独处理。

\subsection*{8.2 Dijkstra with lookahead}

Dijkstra with lookahead 先通过一次从 \(s\) 出发的广度优先搜索计算每个顶点的层数，找出所有瓶颈点，并判断哪些瓶颈点是 marked。这个预处理需要 \(O(m+n)\) 时间。

然后算法运行一个修改版的 Dijkstra：

\begin{itemize}
  \item 非瓶颈点仍然放进堆 \(H\) 中；
  \item 瓶颈点不放进堆，而是放进一个数组 \(B\) 中单独维护；
  \item 输出顺序列表记为 \(L\)。
\end{itemize}

数组 \(B\) 保存尚未加入 \(L\) 的一段瓶颈点，按层数递增排列。初始时，\(B\) 包含从最低层开始的一段瓶颈点，直到第一个 marked 瓶颈点为止；如果没有 marked 瓶颈点，就包含所有瓶颈点。当 \(B\) 中的瓶颈点都被加入 \(L\) 后，再用下一段瓶颈点补充 \(B\)。

扫描一个顶点 \(v\) 时，算法和普通 Dijkstra 类似：处理每条出弧 \(vw\)。如果 \(w\) 尚未被发现，就设置当前距离和父亲；如果 \(w\) 已经 labeled 且距离可以改进，就更新距离和父亲。区别在于：如果 \(w\) 是瓶颈点，不把它插入堆；只有非瓶颈点才进入堆或执行 \texttt{decrease-key}。

算法每轮在两类候选之间选择：

\begin{itemize}
  \item 堆 \(H\) 中当前距离最小的非瓶颈点；
  \item 数组 \(B\) 中当前这段瓶颈点。
\end{itemize}

如果堆中最小的非瓶颈点比 \(B\) 中要处理的瓶颈点更小，就像普通 Dijkstra 一样从堆中删除并扫描它。否则，算法会扫描 \(B\) 中的一段瓶颈点，并把真实距离已经确定、且应该排在当前堆最小元素之前的瓶颈点加入 \(L\)。

这里的“lookahead”指的是：算法会向前看一段瓶颈点，提前计算它们的真实距离，而不是等它们像普通顶点一样通过堆逐个弹出。

\subsection*{8.3 正确性}

Theorem 8.4 说明：Dijkstra with lookahead 是正确的。

证明的关键是维护和普通 Dijkstra 类似的不变式：

\begin{itemize}
  \item 顶点被扫描时，它的当前距离等于真实距离；
  \item 已经放进 \(L\) 的顶点，其当前距离不大于还没放进 \(L\) 的顶点；
  \item 顶点按真实距离非递减顺序加入 \(L\)。
\end{itemize}

算法和普通 Dijkstra 的主要区别是它可能提前扫描瓶颈点。提前扫描之所以安全，是因为瓶颈点支配更高层顶点。在最低层未扫描瓶颈点被处理之前，更高层顶点不可能绕开它而被正确到达。

对于连续的 unmarked 瓶颈点，Lemma 8.3 保证：只要知道第一个瓶颈点的真实距离，后面的瓶颈点可以沿着层数顺序依次得到真实距离。因此，把这一段瓶颈点提前扫描，不会破坏真实距离顺序。

最后还需要证明输出列表 \(L\) 是距离顺序。对每个非源点 \(v\)，设 \(p(v)\) 是算法构造的最短路树中的父亲。证明要说明 \(p(v)\) 一定比 \(v\) 更早加入 \(L\)。如果 \(v\) 或 \(p(v)\) 都是普通情况，这很直接；如果父亲是瓶颈点，算法的 tie-breaking 会保证瓶颈点在相同距离的非瓶颈点之前进入 \(L\)。因此，根据 Lemma 3.1，\(L\) 是距离顺序。

\subsection*{8.4 效率}

效率分析分几部分。

首先，找出并标记瓶颈点需要 \(O(m)\) 时间。每个顶点最多被加入堆或数组一次、扫描一次、再移动到 \(L\) 一次。这些基础操作一共是线性级别。

设瓶颈点数量为 \(b\)。非瓶颈点有 \(n-b\) 个，marked 瓶颈点最多也是 \(n-b\) 个，因为每个 marked 瓶颈点的下一层至少有一个非瓶颈点。由 Lemma 8.1 可知，\(n-b=O(\log D)\)。这说明非瓶颈点数量和 marked 瓶颈点数量都可以用 \(\log D\) 控制。

算法只把非瓶颈点插入堆，所以堆插入次数是 \(n-b\)，因此是 \(O(\log D)\)。类似地，除堆操作和 \(B\) 中搜索之外的额外比较，也可以被 \(O(F-n+1+\log D)\) 控制。

Lemma 8.5 控制 \texttt{delete-min} 的总代价。证明思路是：为了分析，可以假想每个瓶颈点在被扫描前也做一次虚假的插入和立即删除。这样只会增加非瓶颈点的 working-set size，而瓶颈点自己的 working-set size 是 1。于是 Lemma 7.3 的证明仍然适用，可以得到所有 \texttt{delete-min} 的总代价是 \(O(\log D)\)。

Lemma 8.6 控制在数组 \(B\) 中查找的代价。Subcase 2b 中需要在 \(B\) 里找出最大的一段瓶颈点，使它们的当前距离不超过堆中最小非瓶颈点的距离。算法用指数搜索加二分搜索完成。作者把每次搜索移动的瓶颈点集合记为 \(B(v)\)，并利用这些集合能产生许多不同的距离顺序，证明所有搜索总代价也是 \(O(\log D)\)。

最后得到 Theorem 8.7：

\begin{itemize}
  \item Dijkstra with lookahead 的时间复杂度是 \(O(m+\log D)\)；
  \item 在有向图上的比较次数是 \(O(F-n+1+\log D)\)；
  \item 在无向图上的比较次数是 \(O(m-n+1+\log D)\)。
\end{itemize}

因此，Dijkstra with lookahead 在时间和比较次数上都达到了 universal optimality。

\subsection*{这一节的意义}

第七节说明普通 Dijkstra 已经在时间上最优，但比较次数还多一个 \(n\) 级别的项。第八节通过单独处理瓶颈点，把堆中的顶点数和额外比较压到和 \(\log D\) 相关，从而补上了这个差距。

直观地说，如果图里非瓶颈点很多，那么 \(D\) 本来就很大，普通 Dijkstra 已经足够好；如果图里瓶颈点很多，那么 lookahead 会利用瓶颈点的结构，把它们提前、成段地处理掉，避免它们带来不必要的堆比较。

\end{document}
